\documentclass[12pt,a4paper]{article}
\usepackage{ugmath}
\usepackage{float}
\usepackage{placeins}
\usepackage{array}
\usepackage[labelfont=bf,labelsep=space]{caption}
\newcommand{\paren}[1]{\left( #1 \right)}
\newcommand{\alumno}{Ricardo León Martínez}
\newcommand{\materia}{Algebra Lineal I}
\newcommand{\profesor}{Claudia Reynoso Alcántara}
\newcommand{\tarea}{Tarea 10}
\newcommand{\fecha}{27/4/2026}


\begin{document}

    \begin{center}
        {\large\textbf{UNIVERSIDAD DE GUANAJUATO}}\\[0.3cm]
        {\normalsize\textbf{DIVISIÓN DE CIENCIAS NATURALES Y EXACTAS}}\\
        {\normalsize\textbf{CAMPUS GUANAJUATO}}\\[1cm]

        {\Large\textbf{\tarea\ (\materia)}}\\[1cm]
    \end{center}

    \noindent
    \textbf{Nombre:} \alumno 
    \hfill 
    \textbf{Fecha:} \fecha 
    \hfill 
    \textbf{Calificación:} \rule{3cm}{0.4pt}

    \vspace{0.3cm}

    \begin{mdframed}[style=mdpinkbox,frametitle={Ejercicio 1}]
        (2 puntos) Sea $\PP_{3}$ el espacio vectorial sobre $\R$ de polinomios de grado menor
        o igual que 3 con coeficientes en $\R$. Considera la función:
        \begin{align*}
            T:\PP_{3}&\to\R^{4}\\
            f(x)&\mapsto(f(1),f(2),f(3),f(4)).
        \end{align*}
        \begin{enumerate}
            \item[(a)] Demuestra que $T$ es lineal.
            \item[(b)] Usa el teorema de rango y nulidad para demostrar que para todos
            $a_{1},a_{2},a_{3},a_{4}\in\R$, existe $f(x)\in\PP_{3}$ tal que $f(j)=a_{j}$
            para todo $j=1,2,3,4$. 
        \end{enumerate}
    \end{mdframed}

        \begin{proof}
        Para probar que $T$ es lineal, verifiquemos para cualesquiera $f,g\in\PP_{3}$ y $\lambda\in\R$
        se cumple
        \begin{enumerate}
            \item[(a)]
            \begin{enumerate}
                \item[1.] $T(f+g)=T(f)+T(g)$.\\
                En efecto,
                \begin{equation*}
                    T(f+g)=((f+g)(1),(f+g)(2),(f+g)(3),(f+g)(4)).
                \end{equation*}
                Por definición de suma de funciones
                \begin{equation*}
                    (f+g)(x)=f(x)+g(x),
                \end{equation*}
                luego
                \begin{equation*}
                    T(f+g)=(f(1)+g(1),f(2)+g(2),f(3)+g(3),f(4)+g(4))=T(f)+T(g).
                \end{equation*}
                \item[2.] $T(\lambda f)=\lambda T(f)$.\\
                En efecto,
                \begin{equation*}
                    T(\lambda f)=((\lambda f)(1),(\lambda f)(2),(\lambda f)(3), (\lambda f)(4))
                \end{equation*}
                Como $(\lambda f)(x)=\lambda f(x)$, se tiene
                \begin{equation*}
                    T(\lambda f)=(\lambda f(1), \lambda f(2), \lambda f(3), \lambda f(4))=\lambda T(f)
                \end{equation*}
                Así, $T$ es lineal.
            \end{enumerate}
            \item[(b)] Sea $f\in\ker(T)$. Entonces
            \begin{equation*}
                T(f)=(0,0,0,0)
            \end{equation*}
            lo cual implica
            \begin{equation*}
                f(1)=f(2)=f(3)=f(4)=0.
            \end{equation*}
            Es decir, $f$ tiene como raices 1,2,3,4. Pero $f\in\PP_{3}$, por lo que $\deg(f)\leq3$.
            un polinomio de grado $\leq3$ no puede tener 4 raices distintas. Por lo tanto
            \begin{equation*}
                f=0.
            \end{equation*}
            Luego,
            \begin{equation*}
                \ker(T)=\{0\},\qquad\nul(T)=0.
            \end{equation*}
            Así, por el teorema de rango-nulidad,
            \begin{equation*}
                4=0+\rg(T)\implies\rg(T)=4.
            \end{equation*}
            Pero $\Im(T)\subseteq\R^{4}$ y $\dim(\R^{4})=4$. Por lo tanto
            \begin{equation*}
                \Im(T)=\R^{4}.
            \end{equation*}
            La sobreyectividad de $T$ implica que para todo $(a_{1},a_{2},a_{3},a_{4})\in\R^{4}$.
            Existe $f\in\PP_{3}$ tal que
            \begin{equation*}
                T(f)=(a_{1},a_{2},a_{3},a_{4}),
            \end{equation*}
            es decir,
            \begin{equation*}
                f(j)=a_{j}\qquad\text{ para todo }j=1,2,3,4.
            \end{equation*}
        \end{enumerate}
    \end{proof}

    \begin{mdframed}[style=mdpinkbox,frametitle={Ejercicio 2}]
        (2 puntos) Sea $A$ un matriz $m\times n$ con coeficientes en $F$. Demuestra que
        \begin{enumerate}
            \item[(a)] El rango de las filas de $A$ (que por definición es la \textbf{la dimensión
            del subespacio vectorial que general las filas de $A$} vistas como vectores en $F^{n}$)
            mas la dimensión del espacio de soluciones del sistema homogéneo asociado a $A$ es $n$.
            \item[(b)] Si el rango de filas $A$ es $n$ entonces $AX=0$ tiene solo la solución trivial.
            \item[(c)] Si el rango de filas de $A$  es menor que $n$ entonces $AX=0$ tiene una infinidad
            de soluciones.
        \end{enumerate}
    \end{mdframed}

    \begin{mdframed}[style=mdpinkbox,frametitle={Ejercicio 3}]
        (1 punto) Sean $V$ y $W$ espacios vectoriales de dimensión finita sobre $F$ y sea $T:V\to W$ una
        transformación lineal. Demuestra que $T$ es isomorfismo lineal si y solo si, para toda
        $\BB_{V}=\{v_{1},\dots,v_{n}\}$ de $V$ se tiene que $\BB_{W}=\{T(v_{1}),\dots,T(v_{n})\}$
        es base de $W$.
    \end{mdframed}

    \begin{proof}
        Sea $w\in W$. Como $T$ es sobreyectiva, existe $v\in V$ tal que $T(v)=w$. Dado $\BB_{v}$
        es base, existen escalares $a_{1},\dots,a_{n}\in F$ tales que
        \begin{equation*}
            v=\sum_{i=1}^{n}a_{i}v_{i}.
        \end{equation*}
        Aplicando $T$ y usando linealidad
        \begin{equation*}
            w=T(w)=\sum_{i=1}^{n}a_{i}T(v_{i}).
        \end{equation*}
        Luego $\{T(v_{i})\}$ genera $W$. Supongamos
        \begin{equation*}
            \sum_{i=1}^{n}a_{i}T(v_{i})=0.
        \end{equation*}
        Entonces
        \begin{equation*}
            T\paren{\sum_{i=1}^{n}a_{i}v_{i}}=0.
        \end{equation*}
        Como $T$ es inyectiva, $\ker(T)=\{0\}$, por lo que
        \begin{equation*}
            \sum_{i=1}^{n}a_{i}v_{i}=0.
        \end{equation*}
        Dado que $\BB_{v}$ es una base, se sigue que $a_{i}=0$ para todo $i$. Así, $\{T(v_{1}),\dots,T(v_{n})\}$
        es una base de $W$. Supongamos que para toda base $\BB_{V}=\{v_{1},\dots,v_{n}\}$ de $V$, el conjunto
        \begin{equation*}
            \{T(v_{1}),\dots,T(v_{n})\}
        \end{equation*}
        es base de $W$. En particular, fijando una base cualquiera de $V$; Como $\{T(v_{1})\}$ es base
        de $W$, genera $W\implies T$ es sobreyectiva. Ademas, es linealmente independiente, así que si
        \begin{equation*}
            T(v)=0,\qquad v=\sum a_{i}v_{i},
        \end{equation*}
        entonces
        \begin{equation*}
            0=T(v)=\sum a_{i}T(v_{i})\implies a_{i}=0\implies v=0,
        \end{equation*}
        lo cual implica que $\ker(T)=\{0\}$, es decir, $T$ es inyectiva. Por lo tanto,
        $T$ es biyectiva, y en consecuencia un isomorfismo lineal.
    \end{proof}

    \begin{mdframed}[style=mdpinkbox,frametitle={Ejercicio 4}]
        (1 punto) Sea $T:\R^{3}\to\R^{3}$ definido por $(x_{1},x_{2},x_{3},x_{4})\mapsto(3x_{1},x_{1}-x_{2},2x_{1}+x_{2}+x_{3})$.
        ¿Es $T$ isomorfismo lineal? (Justifica tu respuesta). Si es así, encuentra $T^{-1}$.
    \end{mdframed}

    \begin{proof}
        Calculamos $\ker(T)$
        \begin{equation*}
            T(x_{1},x_{2},x_{3})=(0,0,0)
        \end{equation*}
        si y solo si
        \begin{equation*}
            \begin{cases}
                3x_{1}=0,\\
                x_{1}-x_{2}=0,\\
                2x_{1}+x_{2}+x_{3}=0.
            \end{cases}
        \end{equation*}
        Resolvemos: De $3x_{1}=0\implies x_{1}=0$. Entonces $x_{1}-x_{2}=0\implies-x_{2}=0\implies x_{2}=0$.
        Sustituyendo en la tercera se tiene $x_{3}=0$. Por tanto
        \begin{equation*}
            \ker(T)=\{(0,0,0)\}.
        \end{equation*}
        Así, $T$ es inyectiva y asi $T$ es un isomorfismo lineal. Sea $(y_{1},y_{2},y_{3})=T(x_{1},x_{2},x_{3})$,
        es decir
        \begin{equation*}
            \begin{cases}
                y_{1}=3x_{1},\\
                y_{2}=x_{1}-x_{2},\\
                y_{3}=2x_{1}+x_{2}+x_{3}.
            \end{cases}
        \end{equation*}
        Resolvemos: $x_{1}=\frac{y_{1}}{3}$, $x_{2}=\frac{y_{1}}{3}-y_{2}$. Sustituyendo en la tercera
        \begin{equation*}
            y_{3}=y_{1}-y_{2}+x_{3}\implies x_{3}=y_{3}-y_{1}+y_{2}.
        \end{equation*}
        Así,
        \begin{equation*}
            T^{-1}(y_{1},y_{2},y_{3})=(\frac{y_{1}}{3},\frac{y_{1}}{3}-y_{2},y_{3}-y_{1}+y_{2}).
        \end{equation*}
    \end{proof}

    \begin{mdframed}[style=mdpinkbox,frametitle={Ejercicio 5}]
        (2 puntos) Sean $V$ y $W$ espacios vectoriales sobre $F$, de dimensiones $n$ y $m$, respectivamente.
        Demuestra que $V$ y $W$ son isomorfos si y solo si $m=n$.
    \end{mdframed}

    \begin{proof}
        Si $V$ es de dimensión $n$ sobre $F$ entonces al elegir una base ordenada $\BB=\{v_{1},\dots,v_{n}\}$
        la transformación lineal
        \begin{align*}
            \phi_{\BB}:V&\to F^{n}\\
            v&\mapsto[v]_{\BB}
        \end{align*}
        es una biyeccion y por lo tanto es isomorfismo
        \begin{equation*}
            V\cong W.
        \end{equation*}
        Si $\dim_{F}V=\dim_{F}W=n$, entonces
        \begin{equation*}
            V\cong F^{n}\cong W.
        \end{equation*}
        Supongamos que existe un isomorfismo lineal
        \begin{equation*}
            T:V\to W
        \end{equation*}
        Por el teorema de rango-nulidad
        \begin{equation*}
            \nul(T)+\rg(T)=\dim(V).
        \end{equation*}
        Por ser biyección
        \begin{align*}
            \rg(T)&=\dim(W)\\
            \nul(T)&=0
        \end{align*}
        Así,
        \begin{equation*}
            \dim_{f}W=\dim_{F}V.
        \end{equation*}
    \end{proof}

    \begin{mdframed}[style=mdpinkbox,frametitle={Ejercicio 6}]
        (2 puntos) Sea $V$ un espacio vectorial de dimensión finita sobre $F$ y sea $T:V\to V$
        lineal. Demuestra que $T$ es isomorfismo si y solo si para alguna base ordenada $\BB$
        de $V$ se tiene que $[T_{\BB}]$ es una matriz invertible.
    \end{mdframed}

    \begin{proof}
        Supongamos que $T:V\to V$ es un isomorfismo lineal. Entonces existe una transformación
        lineal $S:V\to V$ tal que, para todo $v\in V$,
        \begin{equation*}
            S(T(v))=v,\quad\text{ y }\quad T(S(v))=v.
        \end{equation*}
        Sea $\BB$ una base ordenada de $V$. Entonces, al pasar a matrices
        \begin{equation*}
            [S]_{\BB}[T]_{\BB}=I_{n},\qquad [T]_{\BB}[S]_{\BB}=I_{n}.
        \end{equation*}
        Esto implica que $[T]_{\BB}$ es invertible. Supongamos que existe una base ordenada $\BB$ tal que
        \begin{equation*}
            [T]_{\BB}=A
        \end{equation*}
        es invertible. Sea $A^{-1}$ su inversa. Definimos una aplicación $S:V\to V$ de la siguiente manera.
        Para cada $v\in V$, si $[v]_{\BB}$ es su vector de coordenadas, definimos $S(v)$ como el unico vector
        cuyo vector de coordenadas satisface
        \begin{equation*}
            [S(v)]_{\BB}=A^{-1}[v]_{\BB}.
        \end{equation*}
        Entonces $S$ es lineal y, para todo $v\in V$,
        \begin{align*}
            [S(T(v))]_{\BB}=A^{-1}A[v]_{\BB}=[v]_{\BB},\\
            [T(S(v))]_{\BB}=AA^{-1}[v]_{\BB}=[v]_{\BB}.
        \end{align*}
        Como los vectores tiene las mismas coordenadas en la base $\BB$, se concluye que
        \begin{equation*}
            S(T(v))=v\quad\text{ y }\quad T(S(v))=v\quad \forall v\in V.
        \end{equation*}
        Por lo tanto, $T$ es un isomorfismo lineal.
    \end{proof}

\end{document}